\Needspace{4\baselineskip}
\begin{Example}
\label{ex:integration-operator-fixed-lower-endpoint}
设 $X=C_b([0,L])$, 用下式定义\textbf{积分算子} $\mathcal{I}:X\to X$:
\begin{equation}\label{eq:integration-operator-fixed-lower-endpoint}
\mathcal{I}(f)(x):=\int_0^xf(s)\,\mathrm{d}s,\quad x\in[0,L].
\end{equation}
那么 $\mathcal{N}(\mathcal{I})=\{0\}$, 且 $\mathcal{R}(\mathcal{I})=\{g\in C_b^1([0,L])\mid g(0)=0\}$. 该算子有界; 请读者计算其范数 (Problem~\ref{ex:integration-operator-supremum-norm}).

注意, 也可将此算子看作从 $L^2(0,L)$ 到 $L^2(0,L)$ 的映射. (Hölder 不等式表明, 对每个 $f\in L^2(0,L)$ 及每个 $x\in(0,L)$, $\mathcal{I}(f)(x)$ 都有良好定义.) 再次利用 Hölder 不等式可知该算子有界, 因为\[\begin{aligned}
\|\mathcal{I}(f)\|_{L^2(0,L)}^2&=\int_0^L|\mathcal{I}(f)(x)|^2\,\mathrm{d}x\\
&=\int_0^L\left(\int_0^xf(s)\,\mathrm{d}s\right)^2\,\mathrm{d}x\\
&\leq\int_0^L\left\{\int_0^x\mathrm{d}s\right\}\left\{\int_0^xf(s)^2\,\mathrm{d}s\right\}\,\mathrm{d}x\\
&\leq\dfrac{L^2}{2}\int_0^Lf(s)^2\,\mathrm{d}s\\
&=\dfrac{L^2}{2}\|f\|_{L^2(0,L)}^2.
\end{aligned}\]利用 Theorem~\ref{thm:distribution-primitive-existence}, 可知 $\mathcal{N}(\mathcal{I})=\{0\}$.
\end{Example}

\Needspace{4\baselineskip}
\begin{Example}
\label{ex:differentiation-operator-domain-and-unboundedness}
现在定义一个\textbf{微分算子}. 当然, 只需通过下式定义例如从 $C_b^1([0,L])$ 到 $C_b([0,L])$ 的算子:
\begin{equation}\label{eq:differentiation-operator-pointwise-definition}
Du(x):=\dfrac{\mathrm{d}u}{\mathrm{d}x}(x).
\end{equation}
以这些空间为定义域和值域, 容易证明微分算子有界. 然而, 这样的证明依赖于以下事实: 在定义域上施加了比值域上 ``更强'' 的范数. 在许多情形中, 这样的限制并不令人满意. 因而将微分算子看作从 $L^2(0,L)$ 到 $L^2(0,L)$ 的映射, 其定义域限制为 $\widetilde{\mathcal{D}}(D):=C_b^1([0,L])$, 值域为 $C_b([0,L])$; 即考察\[L^2(0,L)\supset C_b^1([0,L])\ni u\mapsto\dfrac{\mathrm{d}u}{\mathrm{d}x}\in C_b([0,L])\subset L^2(0,L).\]容易看出该算子无界. 只需取
\begin{equation}\label{eq:unbounded-differentiation-exponential-sequence}
u_n(x)=e^{-nx}
\end{equation}
并计算
\begin{equation}\label{eq:unbounded-differentiation-norm-ratio}
\dfrac{\left\|\dfrac{\mathrm{d}u_n}{\mathrm{d}x}\right\|_{L^2(0,L)}}{\|u_n\|_{L^2(0,L)}}=n.
\end{equation}
注意, 在对微分算子求逆时, 无界性并不是真正的问题. (显然要用积分算子实现这一点.) 最重要的困难在于, 上述算子具有非平凡的零空间: 所有常函数的集合. 为解决这一问题, 只需用一个边界条件限制定义域; 例如可考察定义域 $\mathcal{D}(D):=\{u\in C_b^1([0,L])\mid u(0)=0\}$. 在这一定义下, 微分算子为单射, 上面定义的积分算子 $\mathcal{I}$ 是它的逆.

另一个值得关注的问题是, 所选择的算子定义域并不 ``最优''. 遗漏了一些非常显然的候选函数, 例如分片 $C^1$ 函数. 将在下面关于闭算子的部分 (Section~\ref{subsec:closed-operators}) 中讨论这一问题.
\end{Example}
